\scanpage{37}
(\unsure{interprétés|plutôt} comme $n$-poly. rég. orientés munis d'un sommet). \deleted{Enfin, si $\Pi$ est un $n$-polygône combinatoire,} \inserted{la condition d'unimodularité est invariante.}

\uline{NB} Si on a une \deleted{repr.} réal$^{\rm on}$ géom. régulière d'un $n$-polygône comb. $\Pi=(S,\omega)$ ($n$ fini), alors $\psi_r$ sera la \unsure{représ.|repr.} de $\psi_r(\Pi)=(S,\psi_r(\omega))$ définie par la $=$ application $S\longrightarrow E$ -- la condition de régularité sera en effet invariante p.r. au passage de $\omega$ à $\psi_r(\omega)$, \deleted{\unsure{i.e. de}} étant donné que \deleted{le groupe des} l'on a évidemment
\[
\operatorname{Aut}(\Pi)=\operatorname{Aut}(\psi_r(\Pi))\tag*{(27)}
\]
en tant que ss-groupes de $\unsure{\mathfrak S_S|\mathcal S_S|G_S}$, i.e. le normalisateur dans $\mathfrak S_S$ de \inserted{$\omega$ et de} \deleted{\unsure{repr.|$\Psi_r$}} $\psi_r(\omega)$ est évidemment le $=$ (ce qui normalise $\omega$ normalise $\inserted{\omega'=}\psi_r(\omega)$, et inversemt car on a aussi $\omega=\psi_{r'}(\omega')$, si $r'$ est un inverse de $r$ dans $\Psi_n$).

% vertical rule beside the question
\uline{Question} Peut-on avoir \uline{isomorphie} entre un $n$-polygône \inserted{régulier} $P$ et $\psi_r(P)$, ? si $r\in\Psi_n$, $r\neq1$ dans $\Psi_n$ (i.e. si $r\in\mathbb Z_n^{\unsure{\times|*}}$, $r\neq\pm1$ dans $\mathbb Z_n^\times$) ? On trouvera
% archive sheet number: 36
